\documentclass[10pt,letterpaper]{article}
\usepackage{cornell-notes}

\title{Clause 31.06.03.04.02: Get Area Access}
\author{}
\date{}
\setDocTitle{Clause 31.06.03.04.02: Get Area Access}
\setDocSubtitle{C++ 2024}
\setDocOwner{C++ 2024 Cornell Notes}
\setDocDate{\today}
\setCornellCollection{C++ 2024 Cornell Notes}
\setCornellUnitType{Clause}
\setCornellUnitNumber{31.06.03.04.02}
\setCornellUnitTitle{Get Area Access}
\hypersetup{pdftitle={Clause 31.06.03.04.02: Get Area Access},pdfsubject={Cornell notes},pdfauthor={}}

\begin{document}
\maketitle
\begin{CornellOverviewBox}{Overview}
\textbf{Classification:} Standard library - Normative\\[2pt]
\textbf{Learning objective:} Understand the rules, terminology, and semantic consequences associated with Get area access.
\end{CornellOverviewBox}\begin{CornellSummaryBox}{Summary}
{\sffamily\bfseries\color{Accent} Summary}\par\vspace{3pt}Section 31.6.3.4.2 focuses on Get area access. Apply the specified interface together with its mandates, effects, result conditions, exception behavior, and complexity constraints. When applying it, preserve the distinction between mandatory requirements, recommendations, permissions, and behavior deliberately left undefined, unspecified, or implementation-defined.
\end{CornellSummaryBox}

\end{document}
